% GENERATED FILE. Edit canonical lessons, then run npm run generate:books.
% canonical-source-hash: fnv1a64:6b97e7d941ed48f1
% canonical-lessons: KA-C178-bekantale, KA-C178-akasmat, KA-C178-bekadare, KA-C178-yavudakku, KA-C178-andahage

\chapter{Deliberately, By chance, If you want, In case, By the way}
\label{ch:ka-deliberately-by-chance-if-you-want-in-case-by-the-way}
\hlchaptermodality{178}

\begin{chapteropening}
\textbf{By the end of this chapter:} \emph{I can say deliberately, by chance, if you want, in case and by the way.}

Complete the last lesson of chapter 178: I can say deliberately, by chance, if you want, in case and by the way.
\end{chapteropening}

\section[\texorpdfstring{\kn{ಬೇಕಂತಲೇ} (bēkantalē)}{ಬೇಕಂತಲೇ (bēkantalē)}]{\kn{ಬೇಕಂತಲೇ} (bēkantalē) — deliberately, on purpose}

\label{lesson:KA-C178-bekantale}



\begin{quote}
Before the new one: say the Kannada for an effect, then the Kannada for a difference.
\end{quote}

\subsection*{You'll want to know: \kn{ಬೇಕಂತಲೇ}}
\textbf{\kn{ಬೇಕಂತಲೇ}} — \emph{bēkantalē} — \textquotedblleft{}deliberately, on purpose\textquotedblright{}.

\begin{grammarlens}[title={What you've built}]
One more word for this chapter.
\end{grammarlens}

\subsection*{Your turn}
\begin{itemize}
\raggedright
  \item \emph{Say it:} \emph{bēkantalē}
  \item \emph{Say it:} \emph{bēkantalē}, once more
  \item \emph{Say it:} say \emph{bēkantalē}
  \item \emph{Recall:} say the Kannada for an effect, then the Kannada for a difference, then say \emph{bēkantalē} again
\end{itemize}

\begin{tcolorbox}[breakable,colback=teal!4,colframe=teal!35!black,title={Before you move on}]
What does \kn{ಬೇಕಂತಲೇ} mean? (\textquotedblleft{}Deliberately, on purpose\textquotedblright{}.) Say it once more.
\end{tcolorbox}
\section[\texorpdfstring{\kn{ಅಕಸ್ಮಾತ್} (akasmāt)}{ಅಕಸ್ಮಾತ್ (akasmāt)}]{\kn{ಅಕಸ್ಮಾತ್} (akasmāt) — by chance, accidentally}

\label{lesson:KA-C178-akasmat}



\begin{quote}
Before the new one: say the Kannada for a difference, then the Kannada for deliberately.
\end{quote}

\subsection*{You'll want to know: \kn{ಅಕಸ್ಮಾತ್}}
\textbf{\kn{ಅಕಸ್ಮಾತ್}} — \emph{akasmāt} — \textquotedblleft{}by chance, accidentally\textquotedblright{}.

\begin{grammarlens}[title={What you've built}]
One more word for this chapter.
\end{grammarlens}

\subsection*{Your turn}
\begin{itemize}
\raggedright
  \item \emph{Say it:} \emph{akasmāt}
  \item \emph{Say it:} \emph{akasmāt}, once more
  \item \emph{Say it:} say \emph{akasmāt}
  \item \emph{Recall:} say the Kannada for a difference, then the Kannada for deliberately, then say \emph{akasmāt} again
\end{itemize}

\begin{tcolorbox}[breakable,colback=teal!4,colframe=teal!35!black,title={Before you move on}]
What does \kn{ಅಕಸ್ಮಾತ್} mean? (\textquotedblleft{}By chance, accidentally\textquotedblright{}.) Say it once more.
\end{tcolorbox}
\section[\texorpdfstring{\kn{ಬೇಕಾದರೆ} (bēkādare)}{ಬೇಕಾದರೆ (bēkādare)}]{\kn{ಬೇಕಾದರೆ} (bēkādare) — if you want, if necessary}

\label{lesson:KA-C178-bekadare}



\begin{quote}
Before the new one: say the Kannada for deliberately, then the Kannada for by chance.
\end{quote}

\subsection*{You'll want to know: \kn{ಬೇಕಾದರೆ}}
\textbf{\kn{ಬೇಕಾದರೆ}} — \emph{bēkādare} — \textquotedblleft{}if you want, if necessary\textquotedblright{}.

\begin{grammarlens}[title={What you've built}]
One more word for this chapter.
\end{grammarlens}

\subsection*{Your turn}
\begin{itemize}
\raggedright
  \item \emph{Say it:} \emph{bēkādare}
  \item \emph{Say it:} \emph{bēkādare}, once more
  \item \emph{Say it:} say \emph{bēkādare}
  \item \emph{Recall:} say the Kannada for deliberately, then the Kannada for by chance, then say \emph{bēkādare} again
\end{itemize}

\begin{tcolorbox}[breakable,colback=teal!4,colframe=teal!35!black,title={Before you move on}]
What does \kn{ಬೇಕಾದರೆ} mean? (\textquotedblleft{}If you want, if necessary\textquotedblright{}.) Say it once more.
\end{tcolorbox}
\section[\texorpdfstring{\kn{ಯಾವುದಕ್ಕೂ} (yāvudakkū)}{ಯಾವುದಕ್ಕೂ (yāvudakkū)}]{\kn{ಯಾವುದಕ್ಕೂ} (yāvudakkū) — in case, to be on the safe side}

\label{lesson:KA-C178-yavudakku}



\begin{quote}
Before the new one: say the Kannada for by chance, then the Kannada for if you want.
\end{quote}

\subsection*{You'll want to know: \kn{ಯಾವುದಕ್ಕೂ}}
\textbf{\kn{ಯಾವುದಕ್ಕೂ}} — \emph{yāvudakkū} — \textquotedblleft{}in case, to be on the safe side\textquotedblright{}.

\begin{grammarlens}[title={What you've built}]
One more word for this chapter.
\end{grammarlens}

\subsection*{Your turn}
\begin{itemize}
\raggedright
  \item \emph{Say it:} \emph{yāvudakkū}
  \item \emph{Say it:} \emph{yāvudakkū}, once more
  \item \emph{Say it:} say \emph{yāvudakkū}
  \item \emph{Recall:} say the Kannada for by chance, then the Kannada for if you want, then say \emph{yāvudakkū} again
\end{itemize}

\begin{tcolorbox}[breakable,colback=teal!4,colframe=teal!35!black,title={Before you move on}]
What does \kn{ಯಾವುದಕ್ಕೂ} mean? (\textquotedblleft{}In case, to be on the safe side\textquotedblright{}.) Say it once more.
\end{tcolorbox}
\section[\texorpdfstring{\kn{ಅಂದಹಾಗೆ} (andahāge)}{ಅಂದಹಾಗೆ (andahāge)}]{\kn{ಅಂದಹಾಗೆ} (andahāge) — by the way}

\label{lesson:KA-C178-andahage}



\begin{quote}
Before the new one: say the Kannada for if you want, then the Kannada for in case.
\end{quote}

\subsection*{You'll want to know: \kn{ಅಂದಹಾಗೆ}}
\textbf{\kn{ಅಂದಹಾಗೆ}} — \emph{andahāge} — \textquotedblleft{}by the way\textquotedblright{}.

\begin{grammarlens}[title={What you've built}]
One more word for this chapter.
\end{grammarlens}

\subsection*{Your turn}
\begin{itemize}
\raggedright
  \item \emph{Say it:} \emph{andahāge}
  \item \emph{Say it:} \emph{andahāge}, once more
  \item \emph{Say it:} say \emph{andahāge}
  \item \emph{Recall:} say the Kannada for if you want, then the Kannada for in case, then say \emph{andahāge} again
\end{itemize}

\begin{tcolorbox}[breakable,colback=teal!4,colframe=teal!35!black,title={Before you move on}]
What does \kn{ಅಂದಹಾಗೆ} mean? (\textquotedblleft{}By the way\textquotedblright{}.) Say it once more.
\end{tcolorbox}
